\documentclass[12pt]{article}
\usepackage[T1]{fontenc}
\usepackage[utf8]{inputenc}
\usepackage{lmodern}
\usepackage{textcomp}
\usepackage{enumitem}
\usepackage{geometry}
\usepackage{graphicx}
\usepackage{amsmath, amssymb}
\geometry{margin=1in}
\begin{document}
\begin{center}
Chemistry - Graduate Level Assessment 11 \
Total Marks: 40 \
Answer all questions.
\end{center}

\section*{Multiple Choice (10 marks)}
\begin{enumerate}[label=\arabic*.]
\item What is the approximate boiling point at standard pressure of a solution prepared by dissolving 234 g ofNaClin 500 g of H\_2 O?
\begin{enumerate}[label=(\alph*)]
    \item 102.5^\ensuremath{\circC}
    \item 100^\ensuremath{\circC}
    \item 109.5^\ensuremath{\circC}
    \item 106.78^\ensuremath{\circC}
    \item 101.22^\ensuremath{\circC}
\end{enumerate}
\item A substance that is NOT generally considered to be a toxic pollutant in water is
\begin{enumerate}[label=(\alph*)]
    \item nitrogen
    \item a halogenated hydrocarbon
    \item chlorine
    \item copper
    \item arsenic
\end{enumerate}
\item Of the following compounds, which has the lowest melting point?
\begin{enumerate}[label=(\alph*)]
    \item NaCl
    \item LiCl
    \item AlCl 3
    \item KCl
    \item FeCl 3
\end{enumerate}
\item 10.0 grams of urea, CO(NH\_2)\_2, is dissolved in 200 grams of water. What are the boiling and freezing points of the resulting solution?
\begin{enumerate}[label=(\alph*)]
    \item Boiling point is 101.00^\ensuremath{\circC}, Freezing point is 1.55^\ensuremath{\circC}
    \item Boiling point is 101.00^\ensuremath{\circC}, Freezing point is -2.00^\ensuremath{\circC}
    \item Boiling point is 99.57^\ensuremath{\circC}, Freezing point is 1.55^\ensuremath{\circC}
    \item Boiling point is 100.0^\ensuremath{\circC}, Freezing point is 0.0^\ensuremath{\circC}
    \item Boiling point is 100.43^\ensuremath{\circC}, Freezing point is 1.55^\ensuremath{\circC}
\end{enumerate}
\item Calculate the weight in grams of sulfuric acid in 2.00 liters of 0.100 molarsolution. (MW of H\_2 SO\_4 = 98.1.)
\begin{enumerate}[label=(\alph*)]
    \item 10.0 g
    \item 24.4 g
    \item 15.2 g
    \item 22.3 g
    \item 19.6 g
\end{enumerate}
\end{enumerate}

\section*{True / False (10 marks)}
\textit{Answer True or False}
\begin{enumerate}[label=\arabic*.]
\setcounter{enumi}{5}
\item True or False: AlBr 3 requires a Roman numeral in its name due to the variable oxidation state of the cation aluminum.
\item True or False: Diluting 50 ml of 3.50 M H\_2 SO\_4 to achieve a concentration of 2.00 M requires a final volume of 125 ml.
\item True or False: The quotient [\{2 $x^0$\} / \{(2 $x)^0$\}] simplifies to $x^2$, which incorrectly applies the properties of exponents.
\item True or False: To produce 1.50 g of O\_2, 8.45 g of NaNO\_3 must be decomposed according to the balanced reaction.
\item True or False: The zero-point entropy of nitric oxide (NO) is estimated to be 0.00 cal mol^-1 K^-1, indicating no disorder at absolute zero.
\end{enumerate}

\section*{Long Form Response (20 marks)}
\textit{Answer each question with a clear and complete explanation.}
\begin{enumerate}[label=\arabic*.]
\setcounter{enumi}{10}
\item Calculate the cell potential (E) for the electrochemical cell Ag/AgCl/Cl^-(0.0001) and $Ce^4$+(0.01), $Ce^3$+(0.1)/Pt using the provided standard half-cell potentials. Discuss your approach and calculations.
\item Discuss the decomposition of solid mercury(II) oxide (HgO) into oxygen gas (O$_{2}$) and liquid mercury (Hg). Calculate the moles of O$_{2}$ produced from 4.32 g of HgO, considering its molar mass.
\end{enumerate}

\vfill
\noindent\textit{End of Paper}
\end{document}